\section{Finite sequences}

\import{set/finite.tex}
\import{set/finite-ordinal.tex}

Finite sequences are set-coded functions whose domains are von Neumann
natural numbers.  This module develops their constructors and induction
principle; unlike an enumeration, a finite sequence may repeat values.

\begin{proposition}\label{enumeration_is_finite_sequence}
    Suppose $f$ is an enumeration of $X$.
    Then $f$ is a finite sequence in $X$.
\end{proposition}
\begin{proof}
    Take $n$ such that $n$ is a natural number and
        $f\in\bijections{n}{X}$ by \obv{enumeration}.
    $f\in\funs{n}{X}$ by \obv{bijections_to_funs}.
    Follows by \obv{finite_sequence}.
\end{proof}

\begin{proposition}\label{finite_sequence_enumerates_range_iff_injective}
    Suppose $s$ is a finite sequence in $A$.
    Then $s$ is an enumeration of $\ran{s}$ iff $s$ is injective.
\end{proposition}
\begin{proof}
    $\dom{s}$ is a natural number and
        $s\in\funs{\dom{s}}{A}$ by \obv{finite_sequence_domain}.
    $s\in\surj{\dom{s}}{\ran{s}}$ by \obv{fun_to_surj}.
    If $s$ is an enumeration of $\ran{s}$, then $s$ is injective
        by \auto{enumeration,bijections}.
    If $s$ is injective, then
        $s\in\bijections{\dom{s}}{\ran{s}}$
        by \obv{bijections}.
    Follows by \auto{enumeration}.
\end{proof}

\begin{definition}\label{finite_sequences}
    $\finiteSequences{A}=
        \{s\in\pow{\naturals\times A}\mid
            \text{$s$ is a finite sequence in $A$}\}$.
\end{definition}

\begin{proposition}\label{finite_sequences_iff}
    $s\in\finiteSequences{A}$ iff $s$ is a finite sequence in $A$.
\end{proposition}
\begin{proof}
    If $s\in\finiteSequences{A}$, then $s$ is a finite sequence in $A$
        by \obv{finite_sequences}.
    Show that if $s$ is a finite sequence in $A$, then
        $s\in\pow{\naturals\times A}$.
    \begin{subproof}
        Assume $s$ is a finite sequence in $A$.
        $\dom{s}$ is a natural number and
            $s\in\funs{\dom{s}}{A}$ by \obv{finite_sequence_domain}.
        Hence $\dom{s}\subseteq\naturals$
            by \obv{naturals_is_transitive_set,transitiveset_iff_subseteq}.
        $s\in\rels{\dom{s}}{A}$ by \obv{funs_subseteq_rels,subseteq}.
        Hence $s\subseteq\dom{s}\times A$ by \obv{rels_elim}.
        $\dom{s}\times A\subseteq\naturals\times A$
            by \obv{times_subseteq_left}.
        Hence $s\subseteq\naturals\times A$ by \obv{subseteq_transitive}.
        Follows by \obv{pow_iff,subseteq}.
    \end{subproof}
    Follows by \auto{finite_sequences}.
\end{proof}

\begin{definition}\label{sequence_append}
    $\sequenceAppend{s}{x}=\cons{(\dom{s},x)}{s}$.
\end{definition}

\begin{proposition}\label{sequence_append_domain}
    $\dom{\sequenceAppend{s}{x}}=\suc{\dom{s}}$.
\end{proposition}
\begin{proof}
    $\dom{\sequenceAppend{s}{x}}
        =\dom{\cons{(\dom{s},x)}{s}}$ by \obv{sequence_append}.
    Follows by \auto{dom_cons,suc_eq_cons}.
\end{proof}

\begin{proposition}\label{sequence_append_range}
    $\ran{\sequenceAppend{s}{x}}=\cons{x}{\ran{s}}$.
\end{proposition}
\begin{proof}
    $\ran{\sequenceAppend{s}{x}}
        =\ran{\cons{(\dom{s},x)}{s}}$ by \obv{sequence_append}.
    Follows by \auto{ran_cons}.
\end{proof}

\begin{proposition}\label{sequence_append_function}
    Suppose $s$ is a function.
    Then $\sequenceAppend{s}{x}\in
        \funs{\suc{\dom{s}}}{\ran{\sequenceAppend{s}{x}}}$.
\end{proposition}
\begin{proof}
    $s\in\funs{\dom{s}}{\ran{s}}$
        by \obv{function_exact_type_to_range}.
    $\ran{s}\subseteq\cons{x}{\ran{s}}$ by \auto{subseteq,cons_iff}.
    Hence $s\in\funs{\dom{s}}{\cons{x}{\ran{s}}}$
        by \obv{funs_weaken_codom}.
    $\dom{s}\notin\dom{s}$ by \obv{in_irrefl}.
    $x\in\cons{x}{\ran{s}}$ by \obv{cons_left}.
    Hence $\sequenceAppend{s}{x}\in
        \funs{\cons{\dom{s}}{\dom{s}}}{\cons{x}{\ran{s}}}$
        by \auto{sequence_append,extend_function_at_new_point}.
    Hence $\sequenceAppend{s}{x}$ is a function
        by \obv{funs_is_function}.
    Hence $\sequenceAppend{s}{x}\in
        \funs{\dom{\sequenceAppend{s}{x}}}
            {\ran{\sequenceAppend{s}{x}}}$
        by \obv{function_exact_type_to_range}.
    Follows by \obv{sequence_append_domain}.
\end{proof}

\begin{proposition}\label{sequence_append_old_value}
    Suppose $s$ is a function and $i\in\dom{s}$.
    Then $\apply{\sequenceAppend{s}{x}}{i}=s(i)$.
\end{proposition}
\begin{proof}
    $\dom{s}\notin\dom{s}$ by \obv{in_irrefl}.
    $(i,s(i))\in s$ by \obv{function_apply_elim}.
    Hence $(i,s(i))\in\sequenceAppend{s}{x}$
        by \auto{sequence_append,cons_iff}.
    Hence $\sequenceAppend{s}{x}$ is a function
        by \auto{sequence_append_function,funs_is_function}.
    Follows by \obv{function_apply_intro}.
\end{proof}

\begin{proposition}\label{sequence_append_last_value}
    Suppose $s$ is a function.
    Then $\apply{\sequenceAppend{s}{x}}{\dom{s}}=x$.
\end{proposition}
\begin{proof}
    $(\dom{s},x)\in\sequenceAppend{s}{x}$
        by \auto{sequence_append,cons_iff}.
    $\sequenceAppend{s}{x}$ is a function
        by \auto{sequence_append_function,funs_is_function}.
    Follows by \obv{function_apply_intro}.
\end{proof}

\begin{proposition}\label{sequence_append_restriction}
    Suppose $s$ is a function.
    Then $\restrl{\sequenceAppend{s}{x}}{\dom{s}}=s$.
\end{proposition}
\begin{proof}
    $\sequenceAppend{s}{x}$ is a function
        by \auto{sequence_append_function,funs_is_function}.
    $\dom{s}\subseteq\dom{\sequenceAppend{s}{x}}$
        by \auto{sequence_append_domain,subseteq_self_suc_intro}.
    $\restrl{\sequenceAppend{s}{x}}{\dom{s}}$ is a function
        by \obv{restrl_of_function_is_function}.
    $\dom{\restrl{\sequenceAppend{s}{x}}{\dom{s}}}=\dom{s}$
        by \auto{dom_of_restrl_eq_inter,inter_absorb_supseteq_left}.
    Show that for every $i$ we have
        $\apply{\restrl{\sequenceAppend{s}{x}}{\dom{s}}}{i}=s(i)$.
    \begin{subproof}
        Fix $i$.
        \begin{byCase}
        \caseOf{$i\in\dom{s}$.}
            Follows by
                \auto{restrl_of_function_apply,sequence_append_old_value}.
        \caseOf{$i\notin\dom{s}$.}
            $i\notin\dom{\restrl{\sequenceAppend{s}{x}}{\dom{s}}}$.
            Follows by \obv{function_apply_default}.
        \end{byCase}
    \end{subproof}
    Follows by \obv{funext}.
\end{proof}

\begin{proposition}\label{sequence_append_typed}
    Suppose $s\in\funs{n}{A}$ and $x\in A$.
    Then $\sequenceAppend{s}{x}\in\funs{\suc{n}}{A}$.
\end{proposition}
\begin{proof}
    $s$ is a function and $\dom{s}=n$ by \auto{funs_elim}.
    Hence $\sequenceAppend{s}{x}\in
        \funs{\suc{\dom{s}}}{\ran{\sequenceAppend{s}{x}}}$
        by \obv{sequence_append_function}.
    $\ran{s}\subseteq A$ by \obv{funs_ran}.
    Hence $\ran{\sequenceAppend{s}{x}}\subseteq A$
        by \auto{sequence_append_range,cons_iff,subseteq}.
    Hence $\sequenceAppend{s}{x}\in\funs{\suc{\dom{s}}}{A}$
        by \obv{funs_weaken_codom}.
\end{proof}

\begin{corollary}\label{sequence_append_finite_sequence}
    Suppose $s$ is a finite sequence in $A$ and $x\in A$.
    Then $\sequenceAppend{s}{x}$ is a finite sequence in $A$.
\end{corollary}
\begin{proof}
    Follows by
        \auto{finite_sequence_domain,sequence_append_typed,finite_sequence,omega_successor}.
\end{proof}

\begin{proposition}\label{finite_sequence_split_last}
    Suppose $n$ is a natural number and $s\in\funs{\suc{n}}{A}$.
    Then $\restrl{s}{n}\in\funs{n}{A}$ and $s(n)\in A$ and
        $s=\sequenceAppend{\restrl{s}{n}}{s(n)}$.
\end{proposition}
\begin{proof}
    $s$ is a function and $\dom{s}=\suc{n}$ by \auto{funs_elim}.
    $n\subseteq\suc{n}$ by \obv{subseteq_self_suc_intro}.
    Let $r=\restrl{s}{n}$.
    $r$ is a function by \obv{restrl_of_function_is_function}.
    $\dom{r}=n$ by
        \auto{dom_of_restrl_eq_inter,inter_absorb_supseteq_left}.
    For every $i\in\dom{r}$ we have $r(i)=s(i)\in A$
        by \auto{restrl_of_function_apply,funs_type_apply,subseteq}.
    Thus $r$ is a function to $A$ by \auto.
    Hence $r\in\funs{n}{A}$ by \obv{funs_intro}.
    Thus $r$ is a finite sequence in $A$ by \obv{finite_sequence}.
    $n\in\suc{n}$ by \obv{suc_intro_self}.
    Hence $s(n)\in A$ by \obv{funs_type_apply}.
    $\sequenceAppend{r}{s(n)}$ is a function
        by \auto{sequence_append_typed,funs_is_function}.
    $\dom{\sequenceAppend{r}{s(n)}}=\suc{n}$
        by \auto{sequence_append_domain}.
    For every $i\in\suc{n}$ we have
        $\apply{\sequenceAppend{r}{s(n)}}{i}=s(i)$
        by \auto{suc_iff,sequence_append_old_value,sequence_append_last_value,restrl_of_function_apply}.
    $s\subseteq\sequenceAppend{r}{s(n)}$ and
        $\sequenceAppend{r}{s(n)}\subseteq s$
        by \auto{fun_subseteq,subseteq}.
    Hence $s=\sequenceAppend{r}{s(n)}$
        by \obv{subseteq_antisymmetric}.
\end{proof}

\begin{theorem}\label{finite_sequence_induction}
    Suppose $s$ is a finite sequence in $A$.
    Suppose $\emptyset\in C$.
    Suppose for every $r,x$ such that $r\in C$ and $x\in A$ and
        $r$ is a finite sequence in $A$ we have
        $\sequenceAppend{r}{x}\in C$.
    Then $s\in C$.
\end{theorem}
\begin{proof}
    Let $D=\{n\in\naturals\mid
        \text{for every $r\in\funs{n}{A}$ we have $r\in C$}\}$.
    Show that $\emptyset\in D$.
    \begin{subproof}
        $\emptyset\in\naturals$ by \obv{omega_empty}.
        Show that for every $r\in\funs{\emptyset}{A}$ we have $r\in C$.
        \begin{subproof}
            Fix $r\in\funs{\emptyset}{A}$.
            $r\in\rels{\emptyset}{A}$
                by \obv{funs_subseteq_rels,subseteq}.
            Hence $r\subseteq\emptyset\times A$ by \obv{rels_elim}.
            Thus $r=\emptyset$
                by \auto{times_empty_left,subseteq,setext,emptyset}.
        \end{subproof}
    \end{subproof}
    Show that for every $n\in D$ we have $\suc{n}\in D$.
    \begin{subproof}
        Fix $n\in D$.
        Show that for every $r\in\funs{\suc{n}}{A}$ we have $r\in C$.
        \begin{subproof}
            Fix $r\in\funs{\suc{n}}{A}$.
            $\restrl{r}{n}\in\funs{n}{A}$ and
                $r(n)\in A$ and
                $r=\sequenceAppend{\restrl{r}{n}}{r(n)}$
                by \obv{finite_sequence_split_last}.
            $\restrl{r}{n}$ is a finite sequence in $A$
                by \obv{finite_sequence}.
            Hence $\restrl{r}{n}\in C$.
            Thus $\sequenceAppend{\restrl{r}{n}}{r(n)}\in C$.
            Hence $r\in C$.
        \end{subproof}
        $\suc{n}\in\naturals$ by \obv{omega_successor}.
    \end{subproof}
    Hence $D$ is an inductive set by \obv{inductive_set}.
    Thus $\naturals\subseteq D$ by \obv{omega_minimal,subseteq}.
    $\dom{s}$ is a natural number and $s\in\funs{\dom{s}}{A}$
        by \obv{finite_sequence_domain}.
    Hence $\dom{s}\in D$ by \obv{subseteq}.
\end{proof}

\begin{definition}\label{singleton_sequence}
    $\singletonSequence{x}=\{(\emptyset,x)\}$.
\end{definition}

\begin{proposition}\label{singleton_sequence_eq_append}
    $\singletonSequence{x}=\sequenceAppend{\emptyset}{x}$.
\end{proposition}
\begin{proof}
    Follows by
        \obv{singleton_sequence,sequence_append,dom_emptyset,cons_iff,singleton_iff,setext}.
\end{proof}

\begin{proposition}\label{singleton_sequence_typed}
    Suppose $x\in A$.
    Then $\singletonSequence{x}$ is a finite sequence in $A$.
\end{proposition}
\begin{proof}
    $\emptyset$ is a finite sequence in $A$ by \obv{emptyset_finite_sequence}.
    Follows by
        \auto{singleton_sequence_eq_append,sequence_append_finite_sequence}.
\end{proof}
